\textbf{Phần 1. Trắc nghiệm lựa chọn}
\setcounterpageref{ex}{0}

\begin{ex} %[2H11N58]
Đơn vị nào sau đây dùng để đo độ dài lớn nhất?
\choice
{Mét}
{Đề-xi-mét}
{\True Ki-lô-mét}
{Xăng-ti-mét}
\loigiai{
Ki-lô-mét là đơn vị lớn nhất trong các đơn vị trên.\\
\textbf{Đáp án của câu hỏi là C}
}
\end{ex}

\begin{ex} %[2H11N58]
Tờ tiền mệnh giá 1 000 đồng bằng mấy tờ 500 đồng?
\choice
{\True $2$ tờ}
{$1$ tờ}
{$3$ tờ}
{$4$ tờ}
\loigiai{
$500 \times 2 = 1000$ đồng.\\
\textbf{Đáp án của câu hỏi là A}
}
\end{ex}

\begin{ex} %[2H11H58]
Tính: $45\text{ m} + 35\text{ m} = \dots$
\choice
{$70\text{ m}$}
{\True $80\text{ m}$}
{$90\text{ m}$}
{$85\text{ m}$}
\loigiai{
$45 + 35 = 80\text{ m}$.\\
\textbf{Đáp án của câu hỏi là B}
}
\end{ex}

\begin{ex} %[2H11H58]
Tính: $90\text{ km} - 45\text{ km} = \dots$
\choice
{$35\text{ km}$}
{$55\text{ km}$}
{\True $45\text{ km}$}
{$40\text{ km}$}
\loigiai{
$90 - 45 = 45\text{ km}$.\\
\textbf{Đáp án của câu hỏi là C}
}
\end{ex}

\begin{ex} %[2H11H58]
Mai có 2 tờ 200 đồng và 1 tờ 500 đồng. Mai có tất cả số tiền là
\choice
{$800$ đồng}
{$1000$ đồng}
{$700$ đồng}
{\True $900$ đồng}
\loigiai{
$200 \times 2 + 500 = 900$ đồng.\\
\textbf{Đáp án của câu hỏi là D}
}
\end{ex}

\begin{ex} %[2H11H58]
Điền dấu thích hợp: $1\text{ km} \dots 900\text{ m}$.
\choice
{\True $>$}
{$<$}
{$=$}
{Không so sánh được}
\loigiai{
$1\text{ km} = 1000\text{ m} > 900\text{ m}$.\\
\textbf{Đáp án của câu hỏi là A}
}
\end{ex}

\begin{ex} %[2H11V58]
Một cây cầu dài $350\text{ m}$, một cây cầu khác dài $450\text{ m}$. Cả hai cây cầu dài tổng cộng là
\choice
{$700\text{ m}$}
{$900\text{ m}$}
{$1000\text{ m}$}
{\True $800\text{ m}$}
\loigiai{
$350 + 450 = 800\text{ m}$.\\
\textbf{Đáp án của câu hỏi là D}
}
\end{ex}

\begin{ex} %[2H11V58]
Nam có 1 000 đồng, Nam mua 1 cái bánh giá 600 đồng. Số tiền Nam còn lại là
\choice
{$300$ đồng}
{\True $400$ đồng}
{$500$ đồng}
{$200$ đồng}
\loigiai{
$1000 - 600 = 400$ đồng.\\
\textbf{Đáp án của câu hỏi là B}
}
\end{ex}

\begin{ex} %[2H11V58]
Cuộn dây thứ nhất dài $48\text{ m}$, cuộn dây thứ hai ngắn hơn cuộn thứ nhất $15\text{ m}$. Cuộn dây thứ hai dài là
\choice
{$32\text{ m}$}
{$34\text{ m}$}
{\True $33\text{ m}$}
{$35\text{ m}$}
\loigiai{
$48 - 15 = 33\text{ m}$.\\
\textbf{Đáp án của câu hỏi là C}
}
\end{ex}

\begin{ex} %[2H11V58]
Tính giá trị biểu thức: $15\text{ m} \times 2 + 20\text{ m}$.
\choice
{$45\text{ m}$}
{$55\text{ m}$}
{$60\text{ m}$}
{\True $50\text{ m}$}
\loigiai{
$30 + 20 = 50\text{ m}$.\\
\textbf{Đáp án của câu hỏi là D}
}
\end{ex}

\begin{ex} %[2H11V58]
Mỗi chiếc kem giá 200 đồng. Mua 5 chiếc kem như thế hết số tiền là
\choice
{$800$ đồng}
{\True $1\,000$ đồng}
{$700$ đồng}
{$900$ đồng}
\loigiai{
$200 \times 5 = 1000$ đồng.\\
\textbf{Đáp án của câu hỏi là B}
}
\end{ex}

\begin{ex} %[2H11C58]
Hùng có 1 000 đồng, Dũng có 600 đồng. Hùng phải cho Dũng bao nhiêu đồng để số tiền của hai bạn bằng nhau?
\choice
{\True $200$ đồng}
{$400$ đồng}
{$300$ đồng}
{$100$ đồng}
\loigiai{
Hùng hơn Dũng: $1000 - 600 = 400$ đồng.\\
Hùng cho Dũng: $400 : 2 = 200$ đồng.\\
\textbf{Đáp án của câu hỏi là A}
}
\end{ex}

\textbf{Phần 2. Trắc nghiệm đúng sai}
\setcounterpageref{ex}{0}

\begin{ex} %[2H11H58]
Xét tính đúng sai về đại lượng và tiền tệ:
\choiceTF
{\True $1\text{ km} = 1\,000\text{ m}$.}
{\True 2 tờ 500 đồng = 1 000 đồng.}
{$1\text{ m} = 10\text{ cm}$.}
{\True $1\text{ m} = 100\text{ cm}$.}
\loigiai{
\begin{itemize}
    \item $1\text{ km} = 1000\text{ m}$: Đúng.
    \item 2 tờ 500đ = 1000đ: Đúng.
    \item $1\text{ m} = 100\text{ cm} \neq 10\text{ cm}$. Ý 3 Sai, Ý 4 Đúng.
\end{itemize}
\begin{itemize}
    \item a) \textbf{Đúng}.
    \item b) \textbf{Đúng}.
    \item c) \textbf{Sai}.
    \item d) \textbf{Đúng}.
\end{itemize}
}
\end{ex}

\begin{ex} %[2H11V58]
Quãng đường từ nhà đến trường dài $2\text{ km}$. Mỗi ngày An đi học cả đi lẫn về. Xét tính đúng sai:
\choiceTF
{\True Mỗi ngày An đi quãng đường là $4\text{ km}$ ($2 \times 2 = 4$).}
{\True Trong 5 ngày đi học, An đi tất cả $20\text{ km}$ ($4 \times 5 = 20$).}
{Trong 5 ngày An đi tất cả 10 km.}
{\True Nếu chỉ tính 1 chiều đi thì 5 ngày là 10 km, nhưng cả đi và về là 20 km.}
\loigiai{
\begin{itemize}
    \item 1 ngày cả đi về: 4 km: Đúng.
    \item 5 ngày: 20 km: Đúng.
    \item Cả đi và về là 20 km (không phải 10 km). Ý 3 Sai, Ý 4 Đúng.
\end{itemize}
\begin{itemize}
    \item a) \textbf{Đúng}.
    \item b) \textbf{Đúng}.
    \item c) \textbf{Sai}.
    \item d) \textbf{Đúng}.
\end{itemize}
}
\end{ex}

\textbf{Phần 3. Trả lời ngắn}
\setcounterpageref{ex}{0}

\begin{ex} %[2H11H58]
Tính: $36\text{ m} + 24\text{ m}$ theo mét.
\shortans{60}
\loigiai{
\textbf{Đáp số:} $60$\\
$36 + 24 = 60\text{ m}$.
}
\end{ex}

\begin{ex} %[2H11H58]
Tính: $100\text{ km} - 65\text{ km}$ theo ki-lô-mét.
\shortans{35}
\loigiai{
\textbf{Đáp số:} $35$\\
$100 - 65 = 35\text{ km}$.
}
\end{ex}

\begin{ex} %[2H11V58]
Có 3 tờ 200 đồng và 2 tờ 100 đồng. Hỏi có tất cả bao nhiêu đồng?
\shortans{800}
\loigiai{
\textbf{Đáp số:} $800$\\
$600 + 200 = 800$ đồng.
}
\end{ex}

\begin{ex} %[2H11V58]
Lan có 1 000 đồng, mua vở hết 800 đồng. Hỏi Lan còn lại bao nhiêu đồng?
\shortans{200}
\loigiai{
\textbf{Đáp số:} $200$\\
$1000 - 800 = 200$ đồng.
}
\end{ex}

\begin{ex} %[2H11V58]
Sợi dây dài $30\text{ dm}$ cắt thành 5 đoạn bằng nhau. Hỏi mỗi đoạn dài bao nhiêu đề-xi-mét?
\shortans{6}
\loigiai{
\textbf{Đáp số:} $6$\\
$30 : 5 = 6\text{ dm}$.
}
\end{ex}

\begin{ex} %[2H11C58]
Tính nhanh: $450\text{ m} + 280\text{ m} + 250\text{ m} + 20\text{ m}$ theo mét.
\shortans{1000}
\loigiai{
\textbf{Đáp số:} $1000$\\
$(450 + 250) + (280 + 20) = 700 + 300 = 1000\text{ m}$.
}
\end{ex}
